% Diagnóstico de conectividad
\subsection{Diagnóstico de conectividad}\label{sec:3_3_diagnostico_conectividad}
Estas pruebas sirven para confirmar que el firewall funciona y para diagnosticar cuando un servicio no responde. Mientras no se hayan levantado los contenedores (capítulos 5 a 7), los puertos de la plataforma aparecerán cerrados en todos los casos.

\subsubsection{Verificación de puertos y sockets (ss -tulpn, ufw status, iptables -t raw)}
\texttt{ss} muestra qué puertos están escuchando en el servidor y qué proceso los atiende:
\begin{lstlisting}
# ss -tulpn
\end{lstlisting}
\begin{table}[H]
\centering
\caption{Puertos que deben aparecer en escucha con la plataforma instalada}
\small
\begin{tabular}{lp{10.5cm}}
\toprule
\textbf{VPS} & \textbf{Puertos en escucha (\texttt{LISTEN})} \\
\midrule
Frontend & 22 (sshd), 80 y 443 (nginx de router-nginx), 8090--8095 (frontends) \\
Backend & 22 (sshd), 8082--8089, 8091, 8093 y 8094 (microservicios y modelo de IA) \\
Base de datos & 22 (sshd), 5432 (defensoria-db) \\
\bottomrule
\end{tabular}
\end{table}
Que un puerto aparezca en \texttt{LISTEN} solo indica que el servicio está levantado; quién puede llegar a él lo deciden ufw y la regla raw (sección \ref{sec:3_2_reglas_perimetrales}).

ufw registra las conexiones que rechaza con el prefijo \texttt{[UFW BLOCK]}, útil para saber si el firewall está bloqueando algo que debería pasar:
\begin{lstlisting}
# journalctl -k --since "1 hour ago" | grep "UFW BLOCK" | tail
\end{lstlisting}

\subsubsection{Pruebas de conexión entre nodos y desde Internet}
Con \texttt{nc} (incluido en Ubuntu) se prueba si un puerto remoto acepta la conexión; \texttt{-z} solo comprueba sin enviar datos, \texttt{-v} muestra el resultado y \texttt{-w 3} espera como máximo 3 segundos:
\begin{lstlisting}
# nc -zv -w 3 169.58.62.99 5432      # desde el Backend: debe conectar
# nc -zv -w 3 156.67.26.73 8084      # desde el Frontend: debe conectar
\end{lstlisting}
Desde el equipo de administración (fuera de las VPS), los mismos puertos deben fallar por tiempo de espera:
\begin{lstlisting}
$ nc -zv -w 3 169.58.62.99 5432      # debe fallar (timed out)
$ nc -zv -w 3 156.67.26.73 8084      # debe fallar (timed out)
$ nc -zv -w 3 169.58.62.111 443      # debe conectar
\end{lstlisting}

\begin{table}[H]
\centering
\caption{Matriz de pruebas de conectividad (con la plataforma levantada)}
\small
\begin{tabular}{lllc}
\toprule
\textbf{Desde} & \textbf{Hacia} & \textbf{Puerto} & \textbf{Resultado esperado} \\
\midrule
Internet & Frontend 169.58.62.111 & 80 / 443 & Conecta \\
Internet & Backend 156.67.26.73 & 8084 & Falla \\
Internet & BD 169.58.62.99 & 5432 & Falla \\
Frontend & Backend 156.67.26.73 & 8084 & Conecta \\
Frontend & BD 169.58.62.99 & 5432 & Falla \\
Backend & BD 169.58.62.99 & 5432 & Conecta \\
\bottomrule
\end{tabular}
\end{table}
Un resultado \texttt{succeeded} donde se espera \texttt{timed out} indica que falta una regla en esa VPS; uno \texttt{Connection refused} indica que el firewall deja pasar el tráfico pero el contenedor no está levantado.
\captura{red-04-pruebas-nc}{Pruebas de conexión entre nodos y desde Internet}
